% Mathematical sources: the September 29--30 web increment, other_web/03--05,
% and research/A/working_notes/o10_geometry_endpoint_20260929.md, Sections 3--9.
% The preserved originals, rather than this translation, remain the source record.

\section{The fifth determinant power and a positive potential}
\label{sec:web-positive-potential}

All integrals in this section are normalized by the total volume, including
integrals over subsets.  We write
\(\|u\|_{p,\mathrm{av}}=(V^{-1}\int |u|^p)^{1/p}\).
The rough Laplacian is nonnegative, and Hermitian products are conjugate-linear
in their first argument.  Our exterior-algebra normalization is
\(\lvert e^{12}\pm e^{34}\rvert^2=2\).

\begin{theorem}[Finite-density exclusion by a positive line potential]
\label{thm:web-pp-exclusion}
There is no sequence of closed, smooth, actual O10 configurations satisfying
\[
 s_j=-4,\qquad \beta_j\longrightarrow\infty,\qquad
 \kappa_j:=\frac{P_j+D_{\beta_j}}{V_j}\longrightarrow0,
 \qquad \sup_j\fint_{M_j}|W_j^-|^2<\infty.
\]
Consequently, under the actual conformal monopole-supply hypotheses, with
\(p\ge A>0\), \(a^2\ge T_+\), and \(\chi(M)/A\le B\), every fixed
finite \(B\ge0\) admits a constant \(c(B)>0\) such that
\begin{equation}
 \fint_M|\Ric_g+g|^2+\frac{D_\beta(g,a)}{V_g}
 \ge c(B)\frac{(\beta-2)^2}{\beta^2},\qquad \beta\ge2.
 \label{eq:web-pp-finite-density}
\end{equation}
No pointwise curvature bound, noncollapse, diameter bound, or uniform Sobolev
constant is assumed.
\end{theorem}

We give the full positive-potential argument, using the actual-field estimates
of Sections~\ref{sec:web-strong-current} and~\ref{sec:web-direct-resolvent}.
In particular, the proof uses their strong-current conclusions and not merely
a formal spinor with the same pointwise algebraic identities.

\subsection{Actual fields and the remaining derivative}

Suppose that the excluded sequence exists.  Put
\(L=\det S^+\) and \(E=S^-\otimes L\), so that \(\det E=L^3\).
The actual unitary connection on \(L\) is denoted by \(A\); locally write
\(A=i\mathfrak a\), with \(\mathfrak a\) real.  On \(L^k\) write
\(D_k=d+ki\mathfrak a\).  The negative quadratic map has the normalization
\[
 c(q_-(z))z=\frac{|z|^2}{2}z,\qquad
 |q_-(z)|^2=\frac{|z|^4}{8},\qquad Q_z=-iq_-(z).
\]
The actual monopole and its cubic torsion satisfy
\[
 D_A\Phi=0,\qquad F_A^+=2i\sigma(\Phi),\qquad
 \rho=|\Phi|^2\le2,\qquad \rho_*:=\frac43,
\]
\[
 \tau=t_\Phi\in\Gamma(E),\qquad m_*:=\frac{128}{81},\qquad
 |\tau|^2\longrightarrow m_*,\qquad D^-_{3A}\tau\longrightarrow0
 \quad\hbox{in }L^2_{\mathrm{av}}.
\]
The strong-current estimates give uniform integrability of
\(|\nabla\Phi|^4\), a normalized \(L^2\) Hessian bound, and
\begin{equation}
 \rho\longrightarrow\rho_*,\quad
 \zeta:=\langle\Phi,\nabla_A\Phi\rangle\longrightarrow0,\quad
 \nabla\zeta\longrightarrow0,\quad
 N_\tau:=(\nabla\tau)^{1,0}_{J_+}\longrightarrow0
 \quad\hbox{in }L^2_{\mathrm{av}}.
 \label{eq:web-pp-strong-input}
\end{equation}
Moreover, \(\zeta\to0\) in \(L^4_{\mathrm{av}}\).  The bounded
same-bundle resolvent construction supplies \(v\in\Gamma(E)\) with
\begin{equation}
 \begin{gathered}
 |v|+|\nabla^*\nabla v|\le C,\qquad
 v-\tau\longrightarrow0\quad\hbox{in }H^1_{\mathrm{av}}\cap L^4_{\mathrm{av}},\\
 D^-_{3A}v\longrightarrow0,\qquad
 (\nabla^*\nabla-2)v\longrightarrow0,\qquad
 m:=|v|^2\longrightarrow m_*
 \quad\hbox{in }L^2_{\mathrm{av}}.
 \end{gathered}
 \label{eq:web-pp-resolvent-input}
\end{equation}
For \(K:=\sup_j\fint|W_j^-|^2<\infty\), Hessian interchange and
bounded-section interpolation give
\begin{equation}
 \|\nabla v\|_{4,\mathrm{av}}+\|\nabla^2v\|_{2,\mathrm{av}}\le C_K,
 \qquad \fint|\nabla v|^2\longrightarrow2m_*,\qquad
 dm,\ j_v:=\Im\langle v,\nabla v\rangle\longrightarrow0
 \quad\hbox{in }L^2_{\mathrm{av}}.
 \label{eq:web-pp-derivative-input}
\end{equation}

On a fixed region where \(\rho\) and \(m\) are bounded away from zero,
choose adapted frames as below and put
\(Z_1=e_1-ie_2\), \(Z_2=e_3-ie_4\).  The first prolongation gives
\begin{equation}
 \nabla_{Z_1}v,\quad\nabla_{Z_2}v,\quad
 \nabla_{\bar Z_1}v\longrightarrow0
 \quad\hbox{in }L^2_{\mathrm{av}}.
 \label{eq:web-pp-first-prolongation}
\end{equation}
The Dirac trace also makes \((\nabla_{\bar Z_2}v)_1\) small; only
\((\nabla_{\bar Z_2}v)_2\) remains.  Equivalently, the same reduction
follows from the mass, phase, type, and Dirac identities for the bounded
resolvent.  No critical spectral equality or numerical upper bound on \(K\)
is used.  The low-amplitude sets have relative volume tending to zero, and
their \(\nabla v\) energy tends to zero by the \(L^4\) bound in
\eqref{eq:web-pp-derivative-input}.

The actual anti-self-dual curvature fit is
\begin{equation}
 f:=-iF_A^-=-\frac{27}{32}Q_v+o_{L^2_{\mathrm{av}}}(1).
 \label{eq:web-pp-curvature-fit}
\end{equation}
Indeed, the same-field model is
\(K_{\mathrm{fit}}=-(9/8)\rho^3Q_\tau/S\), with \(S\to2m_*\),
and \((9/8)\rho_*^3/(2m_*)=27/32\).  The transfer
\(v-\tau\to0\) in \(L^4\) transfers the quadratic map in \(L^2\).
No derivative of the fitted curvature is taken.

\subsection{The globally defined fifth-power section}

On the good-amplitude region take adapted spinorial and orthonormal frames
such that \(\Phi=\sqrt\rho(1,0)\) and \(v=\sqrt m(1,0)\).  Write
\(\theta_1=e^1-ie^2\), \(\theta_2=e^3-ie^4\).  The connection matrices are
\begin{equation}
 \Omega_+=\begin{pmatrix}ip&-\bar\eta\\\eta&i(\mathfrak a-p)\end{pmatrix},
 \qquad
 \Omega_E=\begin{pmatrix}iq&-\bar b\\b&i(3\mathfrak a-q)\end{pmatrix},
 \label{eq:web-pp-adapted-connections}
\end{equation}
where the real one-forms \(p=\Im\zeta/\rho\) and \(q=j_v/m\) are
phase currents, not the monopole projection square.  The actual cubic
identity and \eqref{eq:web-pp-first-prolongation} yield
\begin{equation}
 \eta=\kappa\theta_1+\epsilon_\eta,
 \quad\kappa=\frac{\sqrt m}{2\rho^{3/2}}\in\mathbb R,
 \qquad b=c_0\theta_2+\epsilon_b,
 \quad c_0=\frac{b_3+ib_4}{2}.
 \label{eq:web-pp-connection-errors}
\end{equation}
Here \(\epsilon_\eta,\epsilon_b,p,q\to0\) in normalized \(L^2\).
All these quantities, as well as \(\eta,b,c_0\), are bounded in normalized
\(L^4\) on every fixed lower-amplitude cutoff.  The bounds are bounds on
covariant derivatives of actual fields, not on the local potential
\(\mathfrak a\).

For a section of \(E\), the invariant cubic expression
\[
 t(v)=\sum_{i=1}^4c^-(e^i)v\otimes(v\wedge\nabla_i v)
 \in\Gamma(S^+\otimes L^4)
\]
satisfies, in these frames,
\[
 t(v)=m^{3/2}(-b_1+ib_2,-b_3-ib_4).
\]
Consequently the coefficient \(c_0\) in
\eqref{eq:web-pp-connection-errors} is the coefficient of the section
\begin{equation}
 c_0=-\frac{\Phi\wedge t(v)}{2\sqrt\rho\,m^{3/2}}
 \quad\hbox{of }L^5.
 \label{eq:web-pp-invariant-section}
\end{equation}
Choose fixed smooth cutoffs \(\chi_\rho,\chi_m\), zero near zero and one
on neighborhoods of \(\rho_*,m_*\), with derivative supports separated
from these limiting amplitudes.  Define globally, by smooth extension by zero,
\begin{equation}
 c=\chi_\rho(\rho)\chi_m(m)c_0\in\Gamma(L^5).
 \label{eq:web-pp-global-section}
\end{equation}
Under a determinant-frame change by \(e^{it}\), the coefficients of
\(c_0,\theta_2,b\) transform respectively by
\(e^{-5it},e^{2it},e^{-3it}\).  Thus the local formula is consistent
with \(D_5=d+5i\mathfrak a\), also directly determined by the wedge in
\eqref{eq:web-pp-invariant-section}.

Differentiating that invariant expression, including its cutoffs, gives
\begin{equation}
 |c|\le C|\nabla v|,
 \qquad |D_5c|\le C\bigl(|\nabla^2v|+|\nabla v|^2
                    +|\nabla\Phi|\,|\nabla v|\bigr).
 \label{eq:web-pp-section-budget}
\end{equation}
Hence \(\|c\|_{4,\mathrm{av}}+\|D_5c\|_{2,\mathrm{av}}\le C_K\).
Its mass is not lost:
\begin{equation}
 \fint|c|^2\longrightarrow1.
 \label{eq:web-pp-positive-mass}
\end{equation}
To see this, on the good-amplitude region
\[
 |\nabla v|^2=\frac{|dm|^2}{4m}+m|q|^2+m|b|^2.
\]
The first two integrals tend to zero.  The connection error gives
\(|b|^2=2|c_0|^2+o_{L^1_{\mathrm{av}}}(1)\); since
\(m-m_*\to0\) in \(L^2\) and \(|b|^2\) is bounded in \(L^2\),
one may replace \(m\) by \(m_*\) in its integral.  The excluded
low-amplitude sets carry vanishing gradient energy.  Now use
\(\fint|\nabla v|^2\to2m_*\).  Boundedness of the fourth moment is
sufficient here; uniform integrability of that fourth moment is not assumed.

\subsection{The Cartan identity and its weak Cauchy--Riemann consequence}

For clarity, the Clifford matrices from \(S^+\) to \(S^-\) are
\[
 C_1=I,\quad C_2=\operatorname{diag}(i,-i),\quad
 C_3=\begin{pmatrix}0&1\\-1&0\end{pmatrix},\quad
 C_4=\begin{pmatrix}0&-i\\-i&0\end{pmatrix};
\]
the reverse matrices are \(-C_i^*\).  Subtracting
\(i\mathfrak a I/2\) and \(3i\mathfrak a I/2\) from the two
matrices in \eqref{eq:web-pp-adapted-connections} gives their traceless
Levi-Civita parts.  The identities
\[
 \Omega_{-,\mathrm{LC}}C_j-C_j\Omega_{+,\mathrm{LC}}
   =\sum_i\Gamma^i{}_jC_i,\qquad
 de^i=-\Gamma^i{}_j\wedge e^j
\]
then give, when \(\eta=\kappa\theta_1\) and \(b=c_0\theta_2\),
\[
 d\theta_2=i(2\mathfrak a-p-q)\wedge\theta_2-\bar c_0\Xi_-,
 \qquad
 \Xi_-=e^{13}+e^{24}+i(e^{14}-e^{23}).
\]
The \(\kappa\) term cancels and \(\Xi_-\) is anti-self-dual.
For the actual connections the same linear computation is
\begin{equation}
 d_{-2A}\theta_2=-i(p+q)\wedge\theta_2-\bar c_0\Xi_-+\mathcal R,
 \qquad |\mathcal R|\le C(|\epsilon_\eta|+|\epsilon_b|).
 \label{eq:web-pp-cartan}
\end{equation}
The off-diagonal curvature of \(E\) is
\((F^E)_{21}=db+i(3\mathfrak a-2q)\wedge b\).  Substitution gives the
full error identity
\begin{equation}
 \begin{split}
 (F^E)_{21}={}&D_5c_0\wedge\theta_2-|c_0|^2\Xi_-
       +d_{3A}\epsilon_b-i(p+3q)\wedge c_0\theta_2\\
       &-2iq\wedge\epsilon_b+c_0\mathcal R.
 \end{split}
 \label{eq:web-pp-offdiagonal-curvature}
\end{equation}
Its self-dual part contains no \(|c_0|^2\Xi_-\) term.  The central
curvature is diagonal, and the remaining self-dual off-diagonal curvature
comes from the trace-free Ricci block.  Thus
\(\lvert((F^E)_{21})^+\rvert\le C|\Ric^\circ|\).  If
\(\mathcal L(\alpha)=(\alpha\wedge\theta_2)^+\), we obtain
\begin{equation}
 \mathcal L(D_5c_0)=-(d_{3A}\epsilon_b)^++G,
 \qquad \|G\|_{4/3,\mathrm{av}}\longrightarrow0.
 \label{eq:web-pp-local-cr}
\end{equation}
For example, the phase terms are products of small \(L^2\) factors and
bounded \(L^4\) factors.  There is no anti-self-dual Weyl term in this
particular self-dual curvature block.

Let \(J_-\) have fundamental form \(\Omega_-=e^{12}-e^{34}\), acting
on one-forms by \((J_-\alpha)(X)=-\alpha(J_-X)\).
In the self-dual basis
\[
 \Psi_1=e^{12}+e^{34},\quad\Psi_2=e^{13}-e^{24},\quad
 \Psi_3=e^{14}+e^{23},\qquad |\Psi_i|^2=2,
\]
the coefficients of \(\mathcal L(\alpha)\) are
\[
 \frac12(-i\alpha_3-\alpha_4,\ \alpha_1+i\alpha_2,
                                    -i\alpha_1+\alpha_2).
\]
With \(\mathcal D=\operatorname{diag}(2,1,1)\), the exact symbol
identity is
\begin{equation}
 \mathcal L^*\mathcal D\mathcal L=I-iJ_-.
 \label{eq:web-pp-cr-symbol}
\end{equation}
The target basis has not been normalized to unit length; its factor of two
is part of this identity.  In particular, the constrained direction is
\(\alpha_1+i\alpha_2=\alpha_3-i\alpha_4=0\), the opposite
Cauchy--Riemann direction to the initial \(J_+\) direction.
On fixed lower-amplitude regions the bundle maps in this computation are
bounded, and their covariant first derivatives are bounded by
\(C(|\nabla\Phi|+|\nabla v|)\).

Introduce the global smooth bounded two-form and its skew endomorphism
\[
 \bar\Omega=-\frac4{m_*}Q_v,\qquad \bar J=J_{\bar\Omega}.
\]
It is not necessary that \(\bar J^2=-I\).  For every sequence of smooth
test sections \(w\in\Gamma(L^5)\) satisfying
\begin{equation}
 \|w\|_\infty+\|D_5w\|_{4,\mathrm{av}}
                  +\|D_5^2w\|_{2,\mathrm{av}}\le C_w,
 \label{eq:web-pp-test-class}
\end{equation}
we claim that
\begin{equation}
 \Re\fint\langle D_5c-i\bar J D_5c,D_5w\rangle\longrightarrow0.
 \label{eq:web-pp-global-weak-cr}
\end{equation}
Here are all the localization errors.  Test the cutoff version of
\eqref{eq:web-pp-local-cr} with \(\mathcal D\mathcal L(D_5w)\).
The \(G\) term tends to zero by the \(L^{4/3}\)--\(L^4\) pairing.
Move the derivative on \(\epsilon_b\) onto this test.  The resulting
adjoint is bounded pointwise by
\[
 C\bigl(|D_5^2w|+
       (|\nabla\Phi|+|\nabla v|+|d\chi|)|D_5w|\bigr),
\]
and hence in normalized \(L^2\).  Its pairing with
\(\epsilon_b\to0\) in \(L^2\) vanishes.  The unitary bundle
connections match in this integration by parts, so no estimate on a bare
gauge potential is needed.  Replacing \(c_0\) with \(c=\chi c_0\)
introduces \(c_0d\chi\).  Both factors have bounded \(L^4\) norms;
the product is bounded in \(L^2\) and supported on sets of vanishing
relative volume.  It therefore tends to zero in \(L^{4/3}\).  Two nested
fixed cutoffs make the maps and tests smooth before either amplitude vanishes.
Finally, on the good region \(Q_v=-m\Omega_-/4\), so
\(\bar J=(m/m_*)J_-\).  Boundedness of \(m\) and its \(L^2\)
convergence imply \(\bar J-J_-\to0\) in \(L^4\).  The remaining
error is controlled by the \(L^4\)--\(L^2\)--\(L^4\) product of
this difference, \(D_5c\), and \(D_5w\).  This proves
\eqref{eq:web-pp-global-weak-cr}.  At no point has a derivative of the
small structural error been asserted to converge strongly.

\subsection{The potential and a legal dual resolvent test}

The Hopf divergence identity, \(D^-_{3A}v\to0\), and \(j_v\to0\)
give \(\delta\bar\Omega\to0\) in \(L^2_{\mathrm{av}}\).
Since \(F^{L^5}=5F_A\), the potential in the line-bundle identity is
\begin{equation}
 \begin{split}
 \mathfrak q&:=-i\langle F^{L^5},\bar\Omega\rangle
       =5\langle f,\bar\Omega\rangle\\
 &\longrightarrow
 5\frac{27}{32}\frac4{m_*}\frac{m_*^2}{8}
       =\frac{135m_*}{64}=\frac{10}{3}
       \quad\hbox{in }L^2_{\mathrm{av}}.
 \end{split}
 \label{eq:web-pp-positive-potential}
\end{equation}
The self-dual curvature is orthogonal to \(\bar\Omega\), and
\(|Q_v|^2=m^2/8\).  This calculation concerns the scalar line \(L^5\);
it is not the Lichnerowicz potential of the rank-two bundle containing
\(t(v)\).  The latter cubic lies in the positive spinor line orthogonal
to \(\Phi\), where the determinant-nine potential tends to \(-7\),
not \(+5\).  Passing to \(\Phi\wedge t(v)\) changes both the bundle
and the operator; the positive value above is computed independently.

For any real global two-form \(\Omega\) and sections \(c,w\) of the
same line bundle, integration by parts and the curvature commutator give
\begin{equation}
 \begin{split}
 &\Re\int\langle D_5c-iJ_\Omega D_5c,D_5w\rangle\\
 &\quad=\Re\int\left\langle c,
   \Delta_{5A}w-i\langle F^{L^5},\Omega\rangle w
                    +i\,\delta\Omega\mathbin{\cdot}D_5w\right\rangle,
 \qquad\Delta_{5A}=D_5^*D_5.
 \end{split}
 \label{eq:web-pp-line-pairing}
\end{equation}
In an orthonormal frame the skew term is
\(i\Omega_{ij}\overline{D_ic}D_jw\); moving \(D_i\) gives
\(-\tfrac i2\Omega_{ij}[D_i,D_j]w\), which fixes the displayed
curvature sign.  The remaining term is the codifferential term.  Neither
closedness of \(\Omega\) nor the identity \(J_\Omega^2=-I\) is required.
Apply this to \(\bar\Omega\).  The potential error is controlled by
\(\|\mathfrak q-10/3\|_2\|c\|_2\|w\|_\infty\), and the drift
error by \(\|\delta\bar\Omega\|_2\|c\|_4\|D_5w\|_4\).
Thus \eqref{eq:web-pp-global-weak-cr} implies
\begin{equation}
 \Re\fint\langle c,(\Delta_{5A}+10/3)w\rangle\longrightarrow0
 \quad\hbox{for every test in \eqref{eq:web-pp-test-class}}.
 \label{eq:web-pp-weak-positive-equation}
\end{equation}

Fix a large \(M\) and a smooth radial cap \(T_M\), equal to the identity
on \(|c|\le M\), with \(|T_M(c)|\le2M\) and
\(\Re\langle c,T_M(c)\rangle\ge0\).  Solve on the same line bundle
\begin{equation}
 (\Delta_{5A}+10/3)w=T_M(c).
 \label{eq:web-pp-dual-resolvent}
\end{equation}
The positive potential gives invertibility and the maximum-principle bound
\(\|w\|_\infty\le3(2M)/10\); its rough Laplacian is also bounded.
We verify rather than assume that this is an admissible test.  For a bounded
section of a unitary bundle, Hessian interchange and divergence interpolation
give, with \(I=\fint|D^2w|^2\), \(J=\fint|Dw|^4\),
\begin{align}
 I&\le C\fint|D^*Dw|^2+C\|w\|_\infty^2\fint|F_D|^2
       +C\fint(|\Ric|+|F_D|)|Dw|^2,\label{eq:web-pp-hessian-test}\\
 J&\le C\|w\|_\infty^2I.\label{eq:web-pp-interpolation-test}
\end{align}
The second inequality follows by integrating the divergence of
\(|Dw|^2\Re\langle w,Dw\rangle\).  In the first inequality all
curvature derivatives have been removed by integration by parts; the tangent
curvature contracts to Ricci.  Cauchy--Schwarz and Young's inequality absorb
the last term using the second inequality.  Here \(F_D=5F_A\), whose
normalized \(L^2\) norm is uniformly bounded by the actual curvature fit;
the Ricci norm is likewise bounded.  There is no \(W^-\) block in this
line connection.  It follows that
\(\|D_5^2w\|_{2,\mathrm{av}}+\|D_5w\|_{4,\mathrm{av}}\le C_M\).
No derivative of \(T_M(c)\) has been used.

Equations~\eqref{eq:web-pp-weak-positive-equation} and
\eqref{eq:web-pp-dual-resolvent} now imply
\(\Re\fint\langle c,T_M(c)\rangle\to0\).  On the other hand,
\[
 \liminf\Re\fint\langle c,T_M(c)\rangle
 \ge1-\limsup\fint_{\{|c|>M\}}|c|^2
 \ge1-\frac{C_K}{M^2}.
\]
Choose \(M\) first, so that the last quantity is positive.  This is the
required contradiction and proves the sequential assertion.

\subsection{Returning to all parameters and all allowed metrics}

If \eqref{eq:web-pp-finite-density} failed at fixed finite \(B\), a
sequence would have its left side divided by
\((\beta-2)^2/\beta^2\) tending to zero.  The case \(\beta=2\) is
vacuous.  The finite-parameter theorem of
Section~\ref{sec:single-parameter} excludes bounded \(\beta\), so
\(\beta\to\infty\) and the defect tends to zero.  The exact conformal
comparison
\[
 \widehat\kappa\le
 \frac{\kappa}{1-\sqrt\kappa-\kappa/16}
\]
passes to the scalar-curvature representative \(s=-4\) while retaining
the actual supply and the topological budget.  On that representative,
\(p/V\le1/(2\pi^2)\), and Gauss--Bonnet yields
\[
 \fint|W^-|^2\le\frac{8\pi^2\chi}{V}-\frac23+\frac{P}{2V}
           \le4B-\frac23+o(1).
\]
This contradicts the sequential assertion and proves the theorem.

\begin{remark}[The remaining dependence on density]
An unrestricted O10 escape sequence, if one exists, must have
\(\fint|W^-|^2\to\infty\).  The endpoint identities
\(p/V\to2/(9\pi^2)\), \(\fint|W^+|^2\to0\), and Gauss--Bonnet
then also force \(\chi/p\to\infty\) and \(\chi/A\to\infty\).
The finite \(K\) budget was used for the original resolvent Hessian bound,
the first-derivative reduction, and the mass and derivative bounds for \(c\).
The absence of \(W^-\) from the final line resolvent does not remove those
earlier uses.  Likewise, this proof gives no O11 sharp coefficient: its
nonzero limiting torsion amplitude is absent at the allowed K\"ahler--Einstein
end, where the relevant smaller energy scale is \(Q=Ah=o(V)\).
\end{remark}

\section{A Chern obstruction with growing Dirichlet energy}
\label{sec:web-growing-energy-chern}

The following quantitative line-bundle lemma does not assume bounded
Dirichlet energy along a sequence.  This is useful independently of the
stronger polynomial density estimate in Section~\ref{sec:web-density-rate}.
Throughout this section, subset integrals are again divided by the total
volume.

\begin{proposition}[Growing-energy logarithmic Chern estimate]
\label{prop:web-gec-lemma}
Let \((M^4,g)\) be closed, smooth, and oriented, and let \(K\to M\)
be a Hermitian line bundle with unitary connection \(\nabla\).  Write
\(F=-iF^K\) for its real curvature.  Suppose
\[
 \fint|\Ric|^2\le C_R,\qquad F=F_0+e_F,
 \qquad |F_0|\le C_F.
\]
The model \(F_0\) need not be closed.  Let \(s\) be a smooth section,
\(|s|\le M_0\), and fix \(\ell>0\).  Put
\[
 \mathsf E=1+\fint|\nabla s|^2,\qquad
 d=\fint(|s|-\ell)^2,\qquad
 \Delta=d+\fint|e_F|^2.
\]
If \(0<\Delta\le e^{-4}\), then
\begin{equation}
 \frac{|c_1(K)^2[M]|}{V}
 \le\frac{C\mathsf E}{\sqrt{\log(1/\Delta)}}.
 \label{eq:web-gec-logarithmic-bound}
\end{equation}
Here \(C\) depends only on \(C_R,C_F,M_0,\ell\), and the dimension,
not on \(\mathsf E\), volume, or any noncollapse or Sobolev constant.
If \(\Delta=0\), the line bundle is trivial and its Chern square vanishes.
In particular, a lower bound \(|c_1(K)^2[M]|/V\ge q_0>0\) implies
\begin{equation}
 \log(1/\Delta)\le C_{q_0}\mathsf E^2.
 \label{eq:web-gec-logarithmic-lower-bound}
\end{equation}
\end{proposition}

\begin{proof}
Let \(L=\nabla^*\nabla\),
\(u_\varepsilon=(I+\varepsilon L)^{-1}s\), and
\(K(\varepsilon)=\fint|Lu_\varepsilon|^2\).  The maximum principle
and spectral calculus give
\begin{equation}
 |u_\varepsilon|\le M_0,\quad
 \fint|u_\varepsilon-s|^2\le\varepsilon\mathsf E,\quad
 \fint|\nabla u_\varepsilon|^2\le\mathsf E.
 \label{eq:web-gec-resolvent-bounds}
\end{equation}
For \(\lambda>0\),
\(\int_0^\infty\lambda^2(1+\varepsilon\lambda)^{-2}
\,d\varepsilon=\lambda\), and the zero eigenvalue contributes zero.
Hence
\[
 \int_0^\infty K(\varepsilon)\,d\varepsilon
       =\fint|\nabla s|^2\le\mathsf E.
\]
Averaging with respect to \(d\varepsilon/\varepsilon\) on
\([\Delta^{1/2},\Delta^{1/4}]\) supplies a parameter satisfying
\begin{equation}
 \Delta^{1/2}\le\varepsilon\le\Delta^{1/4},\qquad
 \varepsilon K(\varepsilon)\le
              \frac{4\mathsf E}{\log(1/\Delta)}.
 \label{eq:web-gec-parameter}
\end{equation}

For any section \(u\) bounded by \(M_0\), the Hessian interchange and
interpolation inequalities
\eqref{eq:web-pp-hessian-test}--\eqref{eq:web-pp-interpolation-test} give
\[
 \int|\nabla^2u|^2\le C\int|Lu|^2+CM_0^2\int|F|^2
             +C\int(|\Ric|+|F|)|\nabla u|^2,
 \qquad
 \int|\nabla u|^4\le CM_0^2\int|\nabla^2u|^2.
\]
Cauchy--Schwarz and Young's inequality therefore give
\begin{equation}
 \fint\bigl(|\nabla^2u|^2+|\nabla u|^4\bigr)
 \le C\left(\fint|Lu|^2+\fint|\Ric|^2+\fint|F|^2\right).
 \label{eq:web-gec-hessian}
\end{equation}
In particular, for the selected \(u=u_\varepsilon\),
\(\fint|\nabla u|^4\le C(K(\varepsilon)+1)\).  The curvature
terms are bounded because \(\|e_F\|_{2,\mathrm{av}}^2\le\Delta\le e^{-4}\).
This estimate neither chooses uniformly controlled coordinate charts nor
uses the full Riemann tensor.

Let \(Z=\{|u|\le3\ell/4\}\).  Since a point of \(Z\) must have
either \(|s|-\ell\) or \(|u-s|\) separated from zero, the amplitude
defect and \eqref{eq:web-gec-resolvent-bounds} imply
\begin{equation}
 \mu(Z):=\frac{\Vol Z}{V}
 \le C(d+\varepsilon\mathsf E)\le C\varepsilon\mathsf E.
 \label{eq:web-gec-low-set}
\end{equation}
The last inequality uses \(d\le\Delta\le\varepsilon\) and
\(\mathsf E\ge1\).  It remains a valid bound even when its right side
exceeds one.

Write \(a=|u|^2\), \(j_u=\Im\langle u,\nabla u\rangle\), and
choose a fixed smooth \(\chi(a)\) equal to zero for \(|u|\le\ell/2\)
and one for \(|u|\ge3\ell/4\).  The real one-form
\(\theta=\chi(a)j_u/a\), extended by zero near the zeros of \(u\),
is globally smooth.  In a local unit frame write
\(u=re^{i\varphi}\) and \(\nabla=d+i\alpha\); then
\(j_u/a=d\varphi+\alpha\).  This proves the exact identity
\[
 d\theta=\chi F+\chi'\,da\wedge j_u/a.
\]
Since the actual curvature \(F\) is closed, Stokes' theorem yields
\begin{equation}
 \left|\frac1V\int F\wedge F\right|
 \le C\fint_Z|F|^2+C\fint_Z|F|\,|\nabla u|^2.
 \label{eq:web-gec-chern-cutoff}
\end{equation}
Split \(F=F_0+e_F\).  The first term is at most
\(C(\mu(Z)+\Delta)\).  The second is bounded by
\(C\sqrt{\mu(Z)(K(\varepsilon)+1)}
 +C\sqrt{\Delta(K(\varepsilon)+1)}\).  Consequently
\[
 \left|\frac1V\int F\wedge F\right|
 \le C\left[\varepsilon\mathsf E+\Delta
   +\sqrt{\varepsilon\mathsf E(K(\varepsilon)+1)}
   +\sqrt{\Delta(K(\varepsilon)+1)}\right].
\]
Set \(\mathsf L=\log(1/\Delta)\ge4\).  Using
\eqref{eq:web-gec-parameter}, this is at most
\[
 C\left[\mathsf E\Delta^{1/4}+\Delta
       +\frac{2\mathsf E}{\sqrt{\mathsf L}}
       +\sqrt{\mathsf E}\Delta^{1/8}
       +\frac{2\sqrt{\mathsf E}\Delta^{1/4}}{\sqrt{\mathsf L}}
       +\sqrt\Delta\right]
 \le\frac{C\mathsf E}{\sqrt{\mathsf L}}.
\]
Finally \(c_1(K)^2[M]=(4\pi^2)^{-1}\int F\wedge F\), which proves
\eqref{eq:web-gec-logarithmic-bound} and its stated consequence.  If
\(\Delta=0\), smoothness and \(|s|=\ell>0\) almost everywhere imply
that \(s\) is a nowhere-zero section, proving the last remaining case.
\end{proof}

\subsection{Quantitative mass and energy for the actual fifth-power section}

We retain the growth of the negative spin connection explicitly.  On an
actual scalar-normalized O10 configuration with \(\beta\ge4\), put
\[
 x=\kappa_\beta+\beta^{-2},\qquad
 b=\fint|W^-|^2,\qquad \mathsf K=1+b.
\]
Use the bounded same-bundle section \(v\) and the strong-current fields
of Sections~\ref{sec:web-strong-current} and~\ref{sec:web-direct-resolvent}.
They obey, without an upper bound on \(b\),
\[
 |\Phi|+|v|+|Tv|\le C,\qquad
 \fint|\nabla\Phi|^4+\fint|\nabla v|^2\le C,
 \qquad T=\nabla_{3A}^*\nabla_{3A}.
\]
Let \(\mathfrak e\) dominate the following primitive errors:
\begin{equation}
 \begin{split}
 &\|\rho-\rho_*\|_2+\||v|^2-m_*\|_2+\|Tv-2v\|_2
       +\|D^-_{3A}v\|_2\\
 &\quad+\|N_v\|_{2,\{\rho\ge c_0\}}
       +\|f-K_{\mathrm{fit}}\|_2+\|\Ric+g\|_2,
 \qquad 0<c_0<\rho_*/4.
 \end{split}
 \label{eq:web-gec-primitive-errors}
\end{equation}
Here all norms are normalized, and \(K_{\mathrm{fit}}\) is the bounded
same-field curvature model.  One may use the larger primitive aggregate
\(\eta_D\) of Proposition~\ref{prop:web-primitive-rate}; fixed comparison
constants are absorbed in \(\mathfrak e\).  For now assume only
\(0<\mathfrak e\le1\).  The complete curvature of \(E=S^-\otimes L\)
has \(\fint|F^E|^2\le C\mathsf K\).  The fixed-rank version of
\eqref{eq:web-gec-hessian} gives
\begin{equation}
 \fint\bigl(|\nabla^2v|^2+|\nabla v|^4\bigr)\le C\mathsf K.
 \label{eq:web-gec-growing-hessian}
\end{equation}
Thus no \(b\)-independent fourth-moment bound is being asserted.

Put \(a=|v|^2\), \(q=|\nabla v|^2\), and \(q_*=2m_*\).
The equation \(Tv=2v+r_v\) yields
\begin{equation}
 \Delta a=4a-2q+r_a,\qquad
 r_a=2\Re\langle v,r_v\rangle,\qquad
 \|r_a\|_2\le C\mathfrak e,\qquad
 \left|\fint q-q_*\right|\le C\mathfrak e.
 \label{eq:web-gec-mass-equation}
\end{equation}
Testing with \(a-m_*\), and using \(\|q\|_2\le C\mathsf K^{1/2}\),
gives
\begin{equation}
 \|da\|_2^2\le C\mathfrak e\mathsf K^{1/2}.
 \label{eq:web-gec-radial}
\end{equation}
For \(\mathcal T(q)=q/(1+q)\),
\(|d\mathcal T(q)|\le C|\nabla^2v|\), so
\[
 \left|\fint\langle da,d\mathcal T(q)\rangle\right|
       \le C\mathfrak e^{1/2}\mathsf K^{3/4}.
\]
Pair \eqref{eq:web-gec-mass-equation} with \(\mathcal T(q)\), and
subtract \(\mathcal T(q_*)\) times its mean identity.  The errors from
\(a-m_*\) and \(r_a\) are at most \(C\mathfrak e\), while
\[
 (\mathcal T(q)-\mathcal T(q_*))(q-q_*)
       =\frac{(q-q_*)^2}{(1+q)(1+q_*)}.
\]
It follows that
\begin{equation}
 \fint\frac{(q-q_*)^2}{1+q}
       \le C\mathfrak e^{1/2}\mathsf K^{3/4},\qquad
 \|q-q_*\|_1\le C\mathfrak e^{1/4}\mathsf K^{3/8}.
 \label{eq:web-gec-covariance}
\end{equation}
The second inequality uses Cauchy--Schwarz and the uniform bound
\(\fint q\le C\), not a growing \(L^1\) budget.

Let \(G=\{\rho\ge c_0,\ a\ge m_*/4\}\).  The amplitude errors and
\eqref{eq:web-gec-growing-hessian} give
\begin{equation}
 \mu(M\setminus G)\le C\mathfrak e^2,\qquad
 \fint_{M\setminus G}q\le C\mathfrak e\mathsf K^{1/2}.
 \label{eq:web-gec-bad-amplitudes}
\end{equation}
The type-projection identity
\(\lvert j_v+J_+da/2\rvert^2\le2a|N_v|^2\), together with
\eqref{eq:web-gec-radial} and \eqref{eq:web-gec-bad-amplitudes}, implies
\(\|j_v\|_2^2\le C\mathfrak e\mathsf K^{1/2}\).
On \(G\), use the adapted frames and write
\[
 A_i=\nabla_{Z_i}v,\qquad B=\nabla_{\bar Z_1}v,\qquad
 C^{\mathrm{vec}}=\nabla_{\bar Z_2}v,\qquad
 c_2=(C^{\mathrm{vec}})_2.
\]
The first components are controlled by
\((\nabla v)_1=(da/2+ij_v)/\sqrt a\), the \(A_i\) by the type
projection, and the remaining unwanted component by the Dirac trace
\[
 D^-v=((A_1)_1-(A_2)_2,\ B_2+(C^{\mathrm{vec}})_1)
\]
up to the immaterial common sign of the reverse Clifford convention.
The total quadratic energy of all components except \(c_2\) is therefore
at most \(C\mathfrak e\mathsf K^{1/2}\).  This reduction uses no
derivative of a bracket identity.

Define the global polynomial section
\(\varsigma=\Phi\wedge t(v)\in\Gamma(L^5)\).
In the adapted frame its norm is
\[
 |\varsigma|^2=\rho a^2|c_2|^2,\qquad
 2q=|A_1|^2+|A_2|^2+|B|^2+|C^{\mathrm{vec}}|^2.
\]
Thus the component estimate controls
\(|\varsigma|^2-2\rho a^2q\) in \(L^1\).  Use
\eqref{eq:web-gec-covariance} for \(q-q_*\), and the \(L^2\) amplitude
errors paired with \(\|q\|_2\le C\mathsf K^{1/2}\) for the coefficient
\(\rho a^2-\rho_*m_*^2\).  On the complement of \(G\), the global
polynomial estimate \(|\varsigma|^2\le Cq\) and
\eqref{eq:web-gec-bad-amplitudes} control the omitted integral.  Hence
\begin{equation}
 \bigl\||\varsigma|^2-\ell_*^2\bigr\|_1
 \le C\bigl(\mathfrak e^{1/4}\mathsf K^{3/8}
                  +\mathfrak e\mathsf K^{1/2}\bigr)
 \le C\mathfrak e^{1/4}\mathsf K,
 \quad
 \ell_*^2=4\rho_*m_*^3=\frac{33554432}{1594323}.
 \label{eq:web-gec-polynomial-mass}
\end{equation}
Direct covariant differentiation, using the bounded amplitudes, gives
\[
 |D_5\varsigma|\le C\bigl(|\nabla\Phi|\,|\nabla v|
                        +|\nabla v|^2+|\nabla^2v|\bigr),
 \qquad 1+\fint|D_5\varsigma|^2\le C\mathsf K.
\]
A fixed smooth radial cap, equal to the identity near \(\ell_*\), gives
a bounded section \(s\in\Gamma(L^5)\) such that
\begin{equation}
 |s|\le M_*,\qquad
 \mathsf E_s:=1+\fint|D_5s|^2\le C\mathsf K,\qquad
 d_s:=\fint(|s|-\ell_*)^2\le C\mathfrak e^{1/4}\mathsf K.
 \label{eq:web-gec-capped-section}
\end{equation}
For the amplitude estimate use
\((r-\ell_*)^2\le|r^2-\ell_*^2|\) and the fixed Lipschitz constant
of the radial cap.

The line curvature \(5(-iF_A)\) is a uniformly bounded model plus an
error of normalized \(L^2\) norm at most \(C\mathfrak e\).  This is a
consequence of the actual determinant-field fit and involves no \(W^-\).
Thus
\begin{equation}
 \Delta_s:=d_s+\|e_{F^{L^5}}\|_2^2
       \le C\mathsf K\mathfrak e^{1/4}.
 \label{eq:web-gec-total-error}
\end{equation}
The same actual curvature identities give
\(\fint|-iF_A^+|^2\to8/9\) and
\(\fint|-iF_A^-|^2\to2/9\), whence
\begin{equation}
 \frac{c_1(L^5)^2[M]}V\longrightarrow\frac{25}{6\pi^2}.
 \label{eq:web-gec-chern-density}
\end{equation}
In particular, for uniformly small primitive error, this density is at least
\(25/(12\pi^2)\).  Proposition~\ref{prop:web-primitive-rate} supplies
a uniform threshold in terms of \(x\).

If \(0<\Delta_s<e^{-4}\), Proposition~\ref{prop:web-gec-lemma} yields
\(\log(1/\Delta_s)\le C\mathsf K^2\).  But
\eqref{eq:web-gec-total-error} gives the inequality in the other direction
\[
 \log(1/\Delta_s)\ge\tfrac14\log(1/\mathfrak e)-\log(C\mathsf K).
\]
If instead \(\Delta_s\ge e^{-4}\), then
\(C\mathsf K\mathfrak e^{1/4}\ge e^{-4}\) gives the same conclusion
directly.  The case \(\Delta_s=0\) is excluded by the positive Chern
density.  We have proved
\begin{equation}
 \log(1/\mathfrak e)\le C(1+b)^2.
 \label{eq:web-gec-primitive-obstruction}
\end{equation}
This argument uses quantitative mass--phase--Dirac reduction, not the weak
Cauchy--Riemann argument of Section~\ref{sec:web-positive-potential}.

\subsection{Exponential escape and finite-density bounds}

\begin{corollary}[An exponential escape constraint]
\label{cor:web-gec-exponential}
Under the actual O10 hypotheses, there are universal positive constants
\(c,C\) such that every scalar-normalized representative with \(\beta\ge4\)
satisfies
\begin{equation}
 \kappa_\beta+\beta^{-2}\ge
       c\exp\bigl[-C(1+\fint|W^-|^2)^2\bigr].
 \label{eq:web-gec-exponential-escape}
\end{equation}
For \(\chi(M)/A\le B\), \(B\ge0\), the original allowed metrics
and all \(\beta\ge2\) satisfy
\begin{equation}
 \fint|\Ric_g+g|^2+\frac{D_\beta(g,a)}{V_g}
 \ge c_1e^{-C_1(1+B)^2}\frac{(\beta-2)^2}{\beta^2}.
 \label{eq:web-gec-exponential-density}
\end{equation}
\end{corollary}

\begin{proof}
Section~\ref{sec:web-primitive-rate} gives, with a fixed
\(0<\epsilon_0\le1\),
\[
 \gamma=\frac{\epsilon_0}{1+\epsilon_0},\qquad
 \vartheta=\frac\gamma{16}>0,\qquad
 \mathfrak e\le Cx^\vartheta.
\]
For small \(x\), combine this with
\eqref{eq:web-gec-primitive-obstruction}:
\(\vartheta\log(1/x)-\log C\le C(1+b)^2\).
Enlarging constants handles all other \(x\), proving
\eqref{eq:web-gec-exponential-escape}.

On a scalar-normalized representative with \(\kappa_\beta\le1\),
the scalar budget and Gauss--Bonnet give
\[
 \frac AV\le\frac1{2\pi^2},\qquad
 b\le4B-\frac23+\frac P{2V}\le C(1+B).
\]
Choose \(m_B=c_0e^{-C_0(1+B)^2}\le1/8\), so that
\(\kappa_\beta+\beta^{-2}\ge m_B\) for \(\beta\ge4\), and let
\(\beta_B=(2/m_B)^{1/2}\ge4\).  For \(\beta\ge\beta_B\),
\(\kappa_\beta\ge m_B/2\).  For \(2\le\beta\le\beta_B\), the
finite-parameter estimate of Section~\ref{sec:single-parameter} gives
\[
 \kappa_\beta\ge c_f\frac{(\beta-2)^2}{\beta^4}
       \ge\frac{c_fm_B}{2}\frac{(\beta-2)^2}{\beta^2}.
\]
If \(\kappa_\beta\ge1\), the required estimate is immediate after
decreasing \(c_0\).  Finally, the exact conformal comparison used in
Section~\ref{sec:web-positive-potential} gives
\(\widehat\kappa\le2\kappa_g\) when \(\kappa_g<1/16\);
larger defects are again harmless.  This transfers the inequality back to
the original allowed metric and proves
\eqref{eq:web-gec-exponential-density}.
\end{proof}

\begin{remark}[Comparison with the later polynomial bound]
Equation~\eqref{eq:web-gec-exponential-escape} excludes, for example,
an escape with \(1+b=o(\sqrt{\log(1/x)})\).  The positive-potential
quantitative continuation, combined with the primitive rate, proves the
stronger estimate \eqref{eq:web-polynomial-density}, namely
\(x\ge c(1+b)^{-24/\gamma}\).  The logarithmic Chern lemma and its
growing-energy proof remain independent useful statements; the exponential
corollary is not the best density dependence now available.  Neither bound
is independent of density, and no optimality of these exponents is claimed.
\end{remark}

\begin{remark}[Why a Dirichlet-energy dependence is necessary]
On a fixed smooth four-manifold take a Hermitian line bundle with nonzero
Chern square and a smooth transverse section \(s_0\).  With a fixed
connection, set
\(s_\varepsilon=s_0/(|s_0|^2+\varepsilon^2)^{1/2}\).
Then \(|s_\varepsilon|\to1\) in \(L^2\), while the connection and
Chern square do not change.  Near the codimension-two zero set, transversality
identifies \(|s_0|\) with the normal radius \(r\); the angular derivative
contributes \(\int_\varepsilon^1r^{-2}r\,dr=O(|\log\varepsilon|)\),
and the remaining contributions are bounded.  Thus the Dirichlet energy
can grow logarithmically while the amplitude defect vanishes.  This is a
line-bundle example, not a counterexample to the actual O10 equations.
\end{remark}

\section{Inverse-amplitude tests without a density bound}
\label{sec:web-inverse-amplitude}

The estimates in this section do not assume a bound on the anti-self-dual
Weyl energy.  They exclude concentration in successively deeper amplitude
wells, but do not assert that all radial concentration disappears.

Let \((M_j,g_j)\) be closed smooth four-manifolds with probability measures
\(d\mu_j=dV_j/V_j\), and let \(E_j\) be Hermitian bundles with unitary
connections.  Write \(L_j=\nabla_j^*\nabla_j\).  Assume that smooth
sections \(v_j\) satisfy
\begin{equation}
 |v_j|+|L_jv_j|\le C,\qquad
 m_j:=|v_j|^2\longrightarrow m_*>0\quad\hbox{in }L^2(d\mu_j),
 \qquad
 r_j:=L_jv_j-\lambda v_j\longrightarrow0\quad\hbox{in }L^2(d\mu_j),
 \label{eq:web-ia-resolvent-hypotheses}
\end{equation}
where \(\lambda>0\) is fixed.  Suppose there are regions \(G_j\) with
orthogonal almost-complex structures \(J_j\), whose orthogonal complex
type projections \(\mathsf P_j=(I-iJ_j)/2\) satisfy
\begin{equation}
 N_j:=\mathsf P_j\nabla_jv_j,
 \qquad \int_{G_j}|N_j|^2\,d\mu_j\longrightarrow0,
 \qquad
 \int_{M_j\setminus G_j}|\nabla_jv_j|^2\,d\mu_j\longrightarrow0.
 \label{eq:web-ia-type-hypotheses}
\end{equation}
Choosing the conjugate type changes none of the argument.
Set \(q_j=|\nabla_jv_j|^2\).  Pairing \(L_jv_j\) with \(v_j\) gives
\(\int q_j\,d\mu_j\to\lambda m_*\).

For the actual O10 section, \(\lambda=2\), \(m_*=128/81\), and one
may take \(G_j=\{\rho_j>c\}\), for any fixed \(0<c<\rho_*\).
The type conclusion is supplied by the strong-current and same-bundle
estimates.  The complementary energy conclusion does not require an
\(L^4\) gradient estimate: choose \(\chi(\rho)=1\) for \(\rho\le c\)
and zero above a second fixed threshold smaller than \(\rho_*\).
Testing \(L_jv_j\) with \(\chi(\rho_j)^2v_j\) and absorbing the cross
term gives
\[
 \int\chi^2|\nabla v_j|^2\,d\mu_j
       \le C\int(\chi^2+|d\chi|^2)\,d\mu_j=o(1).
\]
The last step uses \(\rho_j\to\rho_*\) and \(d\rho_j\to0\) in
\(L^2\), not a commutation of covariant derivatives.

\begin{proposition}[Inverse-amplitude moment]
\label{prop:web-ia-moment}
Under \eqref{eq:web-ia-resolvent-hypotheses} and
\eqref{eq:web-ia-type-hypotheses}, for every fixed \(0<p<1\) and
\(\varepsilon>0\),
\begin{equation}
 \limsup_{j\to\infty}\int
       \frac{q_j}{(m_j+\varepsilon)^p}\,d\mu_j
       \le\frac{\lambda m_*}{(1-p)(m_*+\varepsilon)^p}.
 \label{eq:web-ia-moment}
\end{equation}
No curvature, Hessian, diameter, noncollapse, or Sobolev bound is required.
The parameters \(p,\varepsilon\) are fixed before taking the limit in \(j\).
\end{proposition}

\begin{proof}
On \(G_j\), decompose \(\nabla v=N+A\) into its orthogonal types.
A pure-type complex one-form \(\alpha\) satisfies
\(|2\Re\alpha|^2=2|\alpha|^2\).  Apply this to the scalar one-form
\(\langle v,A\rangle\) and use weighted Young's inequality on the
remaining term.  For every \(\delta>0\),
\begin{equation}
 |dm|^2\le2(1+\delta)m q+C_\delta m|N|^2
       \quad\hbox{on }G_j;
 \qquad |dm|^2\le4mq\quad\hbox{on }M_j\setminus G_j.
 \label{eq:web-ia-refined-kato}
\end{equation}
This is pointwise type algebra, not a regularity assertion.

The smooth global test \((m+\varepsilon)^{-p}v\) in the equation for
\(Lv\) gives exactly
\begin{equation}
 \int\frac{q}{(m+\varepsilon)^p}\,d\mu
 =\Re\int\frac{\langle Lv,v\rangle}{(m+\varepsilon)^p}\,d\mu
       +\frac p2\int\frac{|dm|^2}{(m+\varepsilon)^{p+1}}\,d\mu.
 \label{eq:web-ia-test-identity}
\end{equation}
Call the first term on the right \(F_j\), and the two parts of the left
side \(I_{G,j}\) and \(I_{G^c,j}\).  If \(p(1+\delta)<1\),
\eqref{eq:web-ia-refined-kato} yields
\[
 [1-p(1+\delta)]I_{G,j}
 \le F_j+|2p-1|I_{G^c,j}
       +C_{p,\delta}\varepsilon^{-p}\int_{G_j}|N_j|^2\,d\mu_j.
\]
For fixed \(\varepsilon\), the two error terms tend to zero by
\eqref{eq:web-ia-type-hypotheses}.  Boundedness of \(m_j\), its
\(L^2\) convergence, and the residual convergence show that
\(F_j\to\lambda m_* /(m_*+\varepsilon)^p\).
First take the limit in \(j\), and then let \(\delta\downarrow0\).
This proves \eqref{eq:web-ia-moment}; no sequence-dependent choice of
\(\delta\) or \(\varepsilon\) was made inside a finite-error estimate.
\end{proof}

\begin{corollary}[Energy in deep amplitude wells]
\label{cor:web-ia-deep-wells}
For each fixed \(0<t<m_*\) and \(0<p<1\),
\begin{equation}
 \limsup_j\int_{\{m_j\le t\}}q_j\,d\mu_j
       \le\frac{\lambda m_*}{1-p}\left(\frac t{m_*}\right)^p.
 \label{eq:web-ia-low-amplitude-power}
\end{equation}
For \(0<t<m_*/e\), this implies
\begin{equation}
 \limsup_j\int_{\{m_j\le t\}}q_j\,d\mu_j
       \le e\lambda t\log\frac{m_*}{t},\qquad
 \lim_{t\downarrow0}\limsup_j\int_{\{m_j\le t\}}q_j\,d\mu_j=0.
 \label{eq:web-ia-low-amplitude-log}
\end{equation}
\end{corollary}

\begin{proof}
The integral on \(\{m_j\le t\}\) is bounded by
\((t+\varepsilon)^p\int q_j(m_j+\varepsilon)^{-p}\,d\mu_j\).
Apply Proposition~\ref{prop:web-ia-moment}, then let
\(\varepsilon\downarrow0\).  This proves
\eqref{eq:web-ia-low-amplitude-power}.  For the logarithmic estimate take
the fixed value \(p=1-1/\log(m_*/t)\in(0,1)\) and simplify.
\end{proof}

The order of the two limits in \eqref{eq:web-ia-low-amplitude-log} is
essential.  It does not prove vanishing energy at every fixed
\(0<t<m_*\), convergence \(dm_j\to0\) in \(L^2\), or uniform
integrability of \(q_j\) on arbitrary small-volume sets.  In particular,
it does not establish the fifth-power section's nonzero mass without a
density bound.

This distinction can be seen within the scalar moment identities themselves.
On the amplitude interval \([0,m_*]\), take the formal measures
\[
 d\nu_q(t)=\lambda\,dt,\qquad d\nu_R(t)=2\lambda t\,dt.
\]
For every smooth \(F\), integration by parts gives
\[
 \int F\,d\nu_q+\frac12\int F'\,d\nu_R
       =\lambda m_*F(m_*).
\]
These measures satisfy the limiting mass identity and attain the
inverse-amplitude moment constant as \(\varepsilon\downarrow0\).
They obey deep-well energy vanishing, but
\(\nu_R([0,m_*])=\lambda m_*^2>0\).  Thus the type algebra and the
scalar mass identities used here alone do not exclude energy at fixed
intermediate amplitudes.  These are formal measures, not actual solutions
or a geometric counterexample.

\begin{lemma}[A quadratic-energy relation for bounded line tests]
\label{lem:web-ia-quadratic-test}
Under \eqref{eq:web-ia-resolvent-hypotheses}, let \(w_j\) be smooth
sections of any Hermitian line bundles with unitary rough Laplacians
\(L_w\).  If
\[
 \|w_j\|_\infty+\|L_ww_j\|_\infty
       +\|\nabla w_j\|_{4,\mathrm{av}}\le C,
\]
then
\begin{equation}
 \int(q_j-\lambda m_*)|w_j|^2\,d\mu_j\longrightarrow0.
 \label{eq:web-ia-quadratic-test}
\end{equation}
\end{lemma}

\begin{proof}
The scalar identity is
\(\Delta m_j=2\lambda m_j-2q_j+2\Re\langle r_j,v_j\rangle\).
By self-adjointness of \(\Delta\), the left side of
\eqref{eq:web-ia-quadratic-test} equals
\[
 \lambda\int(m_j-m_*)|w_j|^2\,d\mu_j
 +\int\Re\langle r_j,v_j\rangle|w_j|^2\,d\mu_j
 -\frac12\int(m_j-m_*)\Delta|w_j|^2\,d\mu_j.
\]
The first two terms tend to zero.  The identity
\(\Delta|w_j|^2=2\Re\langle L_ww_j,w_j\rangle
-2|\nabla w_j|^2\) gives a uniform normalized \(L^2\) bound for its
Laplacian, so the third term tends to zero as well.
\end{proof}

The lemma does not require an \(L^2\) bound or uniform integrability for
\(q_j\), but also does not identify \(q_j\) with the fifth-line density
\(|c_j|^2\).  Their intervening radial and phase contributions must still
be controlled.

\section{Low-order geometric transfer of the endpoint form}
\label{sec:web-geometric-transfer}

These results transfer an actual bounded form to an auxiliary metric; they
do not require solving a new monopole equation on that metric.  Their
geometric hypotheses are explicit, and no existence of the auxiliary
regularization is assumed as a consequence of a mean Ricci defect.

On a closed oriented smooth four-manifold let \(*_g\omega=\omega\),
\(|\omega|_g\le L_0\), and \(0\le t\le1/4\).  Define
\begin{equation}
 \begin{split}
 Z_t(g,\omega)={}&|\Ric_g+g|_g^2+
 \left|W_g^+-t\left(\frac{|\omega|_g^2}{3}I-\omega\otimes\omega\right)\right|_g^2\\
 &+(|\omega|_g^2-2)^2+|d\omega|_g^2,
 \qquad z=\frac1{V_g}\int Z_t(g,\omega).
 \end{split}
 \label{eq:web-gt-defect}
\end{equation}
Write
\[
 r_g=\fint_g|\Ric_g+g|^2,\quad
 P_g=\int|\Ric_g^\circ|^2,\quad
 S_g=\int s_g^2,\quad C_g=\int|W_g^+|^2,
\]
where Weyl curvature has its Hilbert--Schmidt norm on two-forms.
The bounded-Ricci--Euler form gap
\eqref{proof:bounded-ricci-endpoint:e3.2} says that
\[
 |\Ric_g|\le K,\quad \frac{8\pi^2\chi(M)}{V_g}\le\Lambda
 \quad\Longrightarrow\quad z\ge c(K,L_0,\Lambda)>0,
\]
including \(t=0\).  That result, with its pointwise Ricci assumption,
is the input to the transfer arguments below.

\subsection{First-derivative budget and self-dual projection}

If \(z\le1\), then
\begin{equation}
 \fint_g|\nabla^g\omega|^2\le C_{L_0}.
 \label{eq:web-gt-form-gradient}
\end{equation}
Indeed, the self-dual Hodge--Weitzenb\"ock identity gives
\[
 \int|\nabla\omega|^2
 =2\int|d\omega|^2-\frac13\int s_g|\omega|^2
                  +2\int\langle W_g^+\omega,\omega\rangle.
\]
The trace estimate is \(\fint_g|s_g|\le4+2\sqrt z\), while the
bounded algebraic model in \eqref{eq:web-gt-defect} gives
\(\fint_g|W_g^+|\le C_{L_0}t+\sqrt z\).  Boundedness of \(\omega\)
now proves \eqref{eq:web-gt-form-gradient}.  No assumption of almost
parallelness or control of \(W^-\) is used.

Let \(h\) be another metric on the same smooth manifold with
\(\delta:=\|h-g\|_{L^\infty(g)}\le\delta_0\), where \(\delta_0>0\)
is a sufficiently small dimensional constant.  The metrics are uniformly
equivalent and \(V_h/V_g=1+O(\delta)\).  Set
\[
 q=\frac1{V_g}\int|\nabla^h-\nabla^g|_g^2\,dV_g,
 \qquad \omega_h=\frac12(\omega+*_h\omega).
\]
Then \(\omega_h\) is \(h\)-self-dual, \(|\omega_h|_h\le CL_0\),
and
\begin{equation}
 \fint_h\bigl[(|\omega_h|_h^2-2)^2+|d\omega_h|_h^2\bigr]
       \le C_{L_0}(z+\delta^2+q).
 \label{eq:web-gt-self-dual-transfer}
\end{equation}
For the proof put \(\mathcal A=(*_h-*_g)/2\).  Since the Hodge stars
are parallel for their respective metric connections,
\[
 |\mathcal A|\le C\delta,\qquad
 |\nabla^g\mathcal A|\le C|\nabla^h-\nabla^g|.
\]
Consequently
\[
 |d\omega_h|_g\le|d\omega|_g+C\delta|\nabla^g\omega|_g
                      +CL_0|\nabla^h-\nabla^g|_g.
\]
Square, integrate, and use \eqref{eq:web-gt-form-gradient}.  The amplitude
estimate follows from
\(\bigl||\omega_h|_h^2-2\bigr|
\le\bigl||\omega|_g^2-2\bigr|+C_{L_0}\delta\).
Uniform metric and volume equivalence prove
\eqref{eq:web-gt-self-dual-transfer}.

\subsection{Topological transfer of the self-dual Weyl energy}

The Chern--Weil identity for the same \(T_+=2\chi+3\tau\) is
\(4\pi^2T_+=2C_g+S_g/24-P_g/2\).  Subtracting its versions for
\(g\) and \(h\) gives the exact formula
\begin{equation}
 C_h-C_g=\frac{S_g-S_h}{48}+\frac{P_h-P_g}{4}.
 \label{eq:web-gt-chern-weil-transfer}
\end{equation}
Let \(r_h=\fint_h|\Ric_h+h|_h^2\).  When \(z,r_h\le1\),
\[
 S_g\le(4+2\sqrt{r_g})^2V_g,\qquad
 S_h\ge(4-2\sqrt{r_h})^2V_h,\qquad P_h\le r_hV_h.
\]
Also \(C_g/V_g\le C_{L_0}t^2+2z\).  Insert these in
\eqref{eq:web-gt-chern-weil-transfer} and discard the favorable
\(-P_g/4\) term to obtain
\begin{equation}
 \frac{C_h}{V_h}\le
       C_{L_0}(t^2+\sqrt z+\delta+\sqrt{r_h}).
 \label{eq:web-gt-weyl-transfer}
\end{equation}
Thus the self-dual Weyl energy transfers to zero when
\(t,z,\delta,r_h\to0\), without estimating the tensor difference
\(W_h^+-W_g^+\) or the change in \(W^-\).  For a nonzero limiting
\(t\), an energy estimate by itself does not retain the prescribed Weyl
model direction; this argument is specifically an endpoint transfer.

\begin{proposition}[A lower bound for regularization cost]
\label{prop:web-gt-cost}
Fix \(K,L_0,\Lambda\).  Suppose \(g,\omega,t\) satisfy
\eqref{eq:web-gt-defect}, \(z\le1\), and
\(8\pi^2\chi/V_g\le\Lambda\).  If \(h\) is a metric on the same
manifold with \(|\Ric_h|_h\le K\) and
\(\|h-g\|_{L^\infty(g)}=\delta\le\delta_0\), then
\begin{equation}
 t^2+\sqrt z+\delta+q+\sqrt{r_h}\ge c(K,L_0,\Lambda)>0.
 \label{eq:web-gt-regularization-cost}
\end{equation}
If additionally \(s_g=s_h=-4\), then, with
\(v_{\mathrm{vol}}=|V_h/V_g-1|\),
\begin{equation}
 t^2+z+v_{\mathrm{vol}}+\delta^2+q+r_h
       \ge c(K,L_0,\Lambda)>0.
 \label{eq:web-gt-normalized-cost}
\end{equation}
\end{proposition}

\begin{proof}
If \(r_h\ge1\), the first inequality is immediate after decreasing
\(c\).  Otherwise \eqref{eq:web-gt-self-dual-transfer} and
\eqref{eq:web-gt-weyl-transfer} imply
\[
 \fint_h Z_0(h,\omega_h)
       \le C_{L_0}(t^2+\sqrt z+\delta+q+\sqrt{r_h}).
\]
The volume comparison implies
\(8\pi^2\chi/V_h\le C\max\{\Lambda,0\}\).  Apply the
bounded-Ricci--Euler form gap to \((h,\omega_h,0)\).
If both scalar curvatures equal \(-4\), then
\(S_g-S_h=16(V_g-V_h)\) in
\eqref{eq:web-gt-chern-weil-transfer}.  Thus its Weyl estimate is instead
\(C_h/V_h\le C_{L_0}(t^2+z+v_{\mathrm{vol}}+r_h)\).
Combine this with \eqref{eq:web-gt-self-dual-transfer} and the same form
gap to prove \eqref{eq:web-gt-normalized-cost}.
\end{proof}

In particular, an endpoint sequence with bounded Euler density,
\(t_j,z_j\to0\), cannot simultaneously admit auxiliary metrics with a
uniform pointwise Ricci bound, \(\|h_j-g_j\|_\infty\to0\), normalized
squared connection difference tending to zero, and \(r_{h_j}\to0\).
For actual O10 fields, the scalar reduction supplies the bounded forms to
which this statement applies.  The finite-density O10 exclusion is already
proved above by a different argument; the transfer proposition remains an
independent geometric restriction and is not an assertion that the needed
auxiliary metrics exist.

\subsection{Connection cost under an actual Bianchi gauge}

There is a conditional improvement when the two metrics have a specified
common identification.  Write \(u=h-g\), and assume under that identification
\begin{equation}
 B_gu:=\delta_gu+\tfrac12d\tr_gu=0.
 \label{eq:web-gt-bianchi-gauge}
\end{equation}
For sufficiently small \(\delta=\|u\|_\infty\),
\begin{equation}
 \int|\nabla^gu|^2\,dV_g
 \le C\delta\int|\Ric_h-\Ric_g|_g\,dV_g
            +C\delta^2\int|\Rm_g|_g\,dV_g.
 \label{eq:web-gt-bianchi-energy}
\end{equation}
Indeed, the principal term of the linearized Ricci difference is
\(\tfrac12\nabla^*\nabla u\).  The other linear terms are the
background curvature acting on \(u\), and the Bianchi divergence term,
which vanishes by \eqref{eq:web-gt-bianchi-gauge}.  The nonlinear terms
have the form \(O(u\nabla^2u)+O(|\nabla u|^2)\), with uniformly bounded
coefficients for small \(u\).  Pair with \(u\) and integrate by parts.
One integration by parts in the first nonlinear term, including the
coefficient derivative, bounds all nonlinear terms by
\(C\delta\int|\nabla u|^2\), which can be absorbed.  The background
curvature term is at most \(C\delta^2\int|\Rm_g|\).  This proves
\eqref{eq:web-gt-bianchi-energy}.

Put \(E_g=\Ric_g+g\) and \(E_h=\Ric_h+h\).  Then
\(\Ric_h-\Ric_g=E_h-E_g-u\).  The full-curvature Gauss--Bonnet
identity and Cauchy--Schwarz give
\[
 \fint_g|\Rm_g|\le C\sqrt{\max\{\Lambda,0\}+r_g}.
\]
The difference of Levi-Civita connections and \(\nabla^gu\) have
uniformly comparable norms.  After dividing
\eqref{eq:web-gt-bianchi-energy} by \(V_g\) and comparing the two
volume forms, we obtain
\begin{equation}
 q\le C\delta(\sqrt{r_g}+\sqrt{r_h})
       +C\delta^2\left(1+\sqrt{\max\{\Lambda,0\}+r_g}\right).
 \label{eq:web-gt-bianchi-cost}
\end{equation}
There is no Sobolev constant or pointwise \(W^-\) bound here.  The
metric gauge \eqref{eq:web-gt-bianchi-gauge} is an additional hypothesis;
this calculation does not produce it.  In particular, it must not be
replaced by an unproved assertion that every smoothing admits a
small-cost common Bianchi identification.  It is also distinct from the
covariant weak field-curvature test of Section~\ref{sec:web-weak-bianchi}.

\subsection{What the classical smoothing statements supply}

The assumptions of the available comparison theorems are stronger than the
mean-normalized hypotheses used here.  In
\href{https://arxiv.org/pdf/1710.09324}{Zerg\"ange, Theorem 1.2}, the
four-manifolds have a bound on the unnormalized total \(L^2\) curvature,
unnormalized trace-free Ricci norm tending to zero, two-sided diameter
bounds, and
\(\Vol B_r(x)\ge\delta\omega_4r^4\) for all small balls, with
\(\delta\) sufficiently close to one.  Those hypotheses do not directly
cover volume-growing or collapsing normalized-defect sequences; its
Gromov--Hausdorff conclusion alone also does not provide the connection
cost above.  \href{https://arxiv.org/pdf/0911.3104}{Li's Ricci-flow smoothing
Theorem 1.1} assumes an initial pointwise Ricci bound and a small full
curvature integral on every ball of the specified radius.  The relevant
four-dimensional estimates of
\href{https://arxiv.org/pdf/1406.6534}{Cheeger--Naber} likewise assume
pointwise Ricci control and noncollapse.  These particular statements do
not construct the auxiliary metrics required here.  This observation does
not exclude a different regularization argument.

\subsection{Replacement inside a set of small capacity}

Global \(C^0\) proximity can be replaced by exact agreement on the
retained region, at the cost of an explicit capacity term.

\begin{proposition}[Small-capacity replacement cost]
\label{prop:web-gt-capacity}
Retain the hypotheses on \(g,\omega,t,z\) and the Euler density from
Proposition~\ref{prop:web-gt-cost}.  Let \(0\le\chi\le1\) be smooth,
put \(\Omega=\{\chi<1\}\), and suppose \(h=g\) on a neighborhood
of \(\{\chi>0\}\).  Inside \(\{\chi=0\}\), the metric \(h\) may
differ arbitrarily, but it is smooth on the same oriented manifold and
satisfies \(|\Ric_h|_h\le K\).  Define
\[
 \alpha_g=\frac{\Vol_g\Omega}{V_g},\qquad
 \alpha_h=\frac{\Vol_h\Omega}{V_g},\qquad
 c_\chi=\frac1{V_g}\int|d\chi|_g^2\,dV_g.
\]
If \(\alpha_g+\alpha_h\le1/2\), then
\begin{equation}
 t^2+\sqrt z+\alpha_g+\sqrt{\alpha_h}+c_\chi
       \ge c(K,L_0,\Lambda)>0.
 \label{eq:web-gt-capacity-cost}
\end{equation}
\end{proposition}

\begin{proof}
The globally smooth form \(\omega_h=\chi\omega\) is \(h\)-self-dual
and bounded by \(L_0\): where it is nonzero, \(h=g\), and at its
zeros the claim is immediate.  All nonzero terms in
\(d(\chi\omega)=d\chi\wedge\omega+\chi d\omega\) are supported where
the metrics agree.  Since \(V_h/V_g\) is bounded above and below,
\[
 \fint_h|d\omega_h|_h^2\le C(z+L_0^2c_\chi),\qquad
 \fint_h(|\omega_h|_h^2-2)^2\le C_{L_0}(z+\alpha_h).
\]
Outside \(\Omega\) the metrics coincide; thus
\[
 \left|\frac{V_h}{V_g}-1\right|\le\alpha_g+\alpha_h,\qquad
 r_h\le C\bigl(z+(K+2)^2\alpha_h\bigr).
\]
Apply the exact identity \eqref{eq:web-gt-chern-weil-transfer} using
this volume comparison, the scalar trace estimates, and the displayed
bound on \(r_h\).  When the proposed left side in
\eqref{eq:web-gt-capacity-cost} is small, \(r_h\le1\) and the previous
calculation gives
\[
 \frac{C_h}{V_h}
       \le C_{K,L_0}(t^2+\sqrt z+\alpha_g+\sqrt{\alpha_h}).
\]
If that left side is not small, the desired lower bound already holds
after decreasing its constant.  It follows in the remaining case that
\(\fint_hZ_0(h,\omega_h)\) is bounded by a fixed multiple of the
left side of \eqref{eq:web-gt-capacity-cost}.  Volume comparability
retains an upper Euler-density bound for \(h\), so the bounded-Ricci
form gap proves the assertion.
\end{proof}

This criterion permits arbitrarily large changes to the metric in the
replaced region and does not require small interior curvature energy,
extension of the monopole or torsion through that region, or transport of
the anti-self-dual Weyl tensor.  It does require both small normalized
volume and small quadratic cutoff capacity, and an actual bounded-Ricci
filling agreeing with the retained exterior.  Small volume alone implies
neither of these latter requirements.  No such replacement construction
is asserted from the weak mean Ricci defect.
